% ============================================================
% 02_DADOS.TEX
% Sistema Adaptativo de Classificação Multicritério
% de Fluidos de Corte
%
% FINALIDADE
%
% Apresentar:
% - origem dos dados;
% - estrutura da base;
% - unidades de análise;
% - condições experimentais;
% - tipos de informação disponíveis.
%
% Este arquivo constitui apenas a BASE da apresentação.
%
% Os valores quantitativos deverão ser preenchidos somente
% depois da auditoria conduzida por Pacheco.
%
% ============================================================


% ============================================================
% FRAME 1 — FONTE DOS DADOS
% ============================================================

\begin{frame}{Fonte dos dados}

    \begin{block}{Base fornecida para o projeto}
        Os dados utilizados no estudo são provenientes do material
        experimental disponibilizado para o trabalho.
    \end{block}

    \vspace{0.3cm}

    O arquivo original será tratado como:

    \begin{itemize}
        \item fonte primária dos dados;
        \item arquivo imutável durante a análise;
        \item referência para auditoria e rastreabilidade;
        \item ponto de partida para a construção da base canônica.
    \end{itemize}

    \vspace{0.3cm}

    \begin{alertblock}{Regra de integridade}
        Nenhuma correção deverá ser realizada diretamente no arquivo
        original. Eventuais ajustes serão realizados de forma
        reproduzível e documentada em R.
    \end{alertblock}

\end{frame}


% ============================================================
% FRAME 2 — ESTRUTURA DO ARQUIVO
% ============================================================

\begin{frame}{Estrutura dos dados}

    \begin{columns}[T]

        \begin{column}{0.48\textwidth}

            \begin{block}{Arquivo original}

                A auditoria deverá identificar:

                \begin{itemize}
                    \item número de abas;
                    \item função de cada aba;
                    \item variáveis disponíveis;
                    \item tipos de dados;
                    \item unidades;
                    \item fórmulas;
                    \item campos ausentes.
                \end{itemize}

            \end{block}

        \end{column}


        \begin{column}{0.48\textwidth}

            \begin{block}{Elementos de análise}

                A estrutura poderá conter informações associadas a:

                \begin{itemize}
                    \item fluidos avaliados;
                    \item condições experimentais;
                    \item desempenho de usinagem;
                    \item qualidade;
                    \item propriedades técnicas;
                    \item informações econômicas;
                    \item especificações.
                \end{itemize}

            \end{block}

        \end{column}

    \end{columns}

    \vspace{0.3cm}

    \centering
    \textit{A composição definitiva da base será confirmada somente após
    a auditoria.}

\end{frame}


% ============================================================
% FRAME 3 — UNIDADE DE ANÁLISE
% ============================================================

\begin{frame}{Unidade de análise}

    A unidade central de decisão será cada alternativa de fluido
    identificada na base validada.

    \vspace{0.3cm}

    Entretanto, os valores associados a cada fluido podem depender de:

    \begin{itemize}
        \item diferentes condições experimentais;
        \item diferentes medições;
        \item repetições;
        \item variáveis com escalas distintas;
        \item informações calculadas ou agregadas.
    \end{itemize}

    \vspace{0.3cm}

    \begin{block}{Implicação}
        Antes da classificação multicritério será necessário estabelecer
        claramente como cada observação será transformada em informação
        comparável entre os fluidos.
    \end{block}

\end{frame}


% ============================================================
% FRAME 4 — CONDIÇÕES EXPERIMENTAIS
% ============================================================

\begin{frame}{Condições experimentais}

    \begin{block}{Preservação da estrutura experimental}
        Quando uma mesma característica for observada em diferentes
        condições, essas condições serão inicialmente analisadas
        separadamente.
    \end{block}

    \vspace{0.3cm}

    Isso permite investigar:

    \begin{itemize}
        \item estabilidade do fluido;
        \item variação entre condições;
        \item possíveis inversões de desempenho;
        \item existência de trade-offs;
        \item adequação de uma eventual agregação.
    \end{itemize}

    \vspace{0.3cm}

    \begin{alertblock}{Evitar agregação prematura}
        Médias ou outros resumos não deverão ser utilizados
        automaticamente antes de verificar se escondem diferenças
        relevantes entre condições.
    \end{alertblock}

\end{frame}


% ============================================================
% FRAME 5 — TIPOS DE INFORMAÇÃO
% ============================================================

\begin{frame}{Tipos de informação}

    Os dados poderão exercer funções diferentes dentro do sistema
    multicritério.

    \vspace{0.2cm}

    \begin{columns}[T]

        \begin{column}{0.48\textwidth}

            \begin{block}{Desempenho}
                Variáveis potencialmente utilizadas para comparar
                quantitativamente as alternativas.
            \end{block}

            \begin{block}{Conformidade}
                Informações associadas a limites ou requisitos técnicos.
            \end{block}

        \end{column}


        \begin{column}{0.48\textwidth}

            \begin{block}{Diagnóstico}
                Variáveis úteis à interpretação, mas que podem não
                participar diretamente do ranking.
            \end{block}

            \begin{block}{Informação auxiliar}
                Dados utilizados para auditoria, cálculo ou contextualização.
            \end{block}

        \end{column}

    \end{columns}

\end{frame}


% ============================================================
% FRAME 6 — DA PLANILHA À BASE CANÔNICA
% ============================================================

\begin{frame}{Da planilha à base analítica}

    \begin{block}{Fluxo previsto}
        \centering

        Arquivo original
        $\rightarrow$
        Inventário
        $\rightarrow$
        Auditoria
        $\rightarrow$
        Reconciliação
        $\rightarrow$
        Base canônica
    \end{block}

    \vspace{0.4cm}

    A base canônica será a versão utilizada pelas etapas estatísticas
    posteriores.

    \vspace{0.3cm}

    Ela deverá possuir:

    \begin{itemize}
        \item estrutura conhecida;
        \item unidades documentadas;
        \item identificadores consistentes;
        \item valores rastreáveis;
        \item problemas pendentes registrados;
        \item transformações reproduzíveis.
    \end{itemize}

\end{frame}


% ============================================================
% FRAME 7 — INFORMAÇÕES A PREENCHER DEPOIS
% ============================================================

\begin{frame}{Caracterização da base}

    \begin{block}{Informações que serão obtidas na auditoria}

        \begin{itemize}
            \item Número de alternativas:
                  \texttt{TO\_BE\_FILLED}

            \item Número de abas:
                  \texttt{TO\_BE\_FILLED}

            \item Número de variáveis relevantes:
                  \texttt{TO\_BE\_FILLED}

            \item Condições experimentais:
                  \texttt{TO\_BE\_FILLED}

            \item Base canônica:
                  \texttt{TO\_BE\_FILLED}
        \end{itemize}

    \end{block}

    \vspace{0.3cm}

    \centering
    \textit{Esses valores deverão ser inseridos pelos integrantes após
    a execução e validação da auditoria em R.}

\end{frame}


% ============================================================
% FIGURA FUTURA
% ============================================================
%
% Caso seja útil, os integrantes poderão substituir parte desta
% seção por uma figura produzida a partir da auditoria.
%
% Exemplo conceitual:
%
% arquivo original
%     ↓
% abas
%     ↓
% variáveis
%     ↓
% base canônica
%
% Não é necessário criar essa figura agora.
%
% ============================================================


% ============================================================
% MATERIAL EXTERNO DA PROFESSORA
% ============================================================
%
% Caso a professora forneça:
%
% - descrição oficial dos dados;
% - significado de abas;
% - definição das variáveis;
% - especificações;
% - unidades;
% - condições experimentais;
%
% utilizar o conteúdo fornecido por ela como referência oficial.
%
% Não preencher lacunas com suposições.
%
% ============================================================


% ============================================================
% ESTADO ATUAL
% ============================================================
%
% DATA_SECTION_STATUS = BASE_PREPARED
%
% N_FLUIDS = TO_BE_FILLED
% N_SHEETS = TO_BE_FILLED
% N_VARIABLES = TO_BE_FILLED
% EXPERIMENTAL_CONDITIONS = TO_BE_FILLED
%
% ============================================================